%% kapa: exam=1 round=36 year=2025 period=1 form=A subject=civil_law subject_ko=민법 seq=14 number=14
%% source: Q-Net 감정평가사 기출문제 (cst00309) artlSeq=5231525 file=raw/1008/5231525/2025년 제36회 감정평가사 1차시험 1교시 문제지 원본.pdf
%% conversion: kordoc 4.24.51 -> tools/gen1.py (auto); review=proofread-sonnet; flags=none
\begin{question}{14}
무권대리행위의 추인에 관한 설명으로 옳지 않은 것은? (다툼이 있으면 판례에 따름)
\begin{choices}
  \item 무권대리행위의 추인은 묵시적인 방법으로도 할 수 있다.
  \item 상대방이 유효하게 무권대리행위를 철회한 후에도 본인은 추인할 수 있다.
  \item 본인이 무권대리인에게 무권대리행위를 추인한 경우에 상대방이 이를 알지 못하는 동안에는 본인은 상대방에게 추인의 효과를 주장하지 못한다.
  \item 무권대리행위의 추인은 상대방의 동의나 승낙을 요하지 않는다.
  \item 무권대리행위의 일부에 대한 추인은 상대방의 동의를 얻지 못하는 한 무효이다.
\end{choices}
\end{question}
\answer{2}
